% !TEX root = ../../main.tex

\section{Pipeline de Validación}

Esta sección de anexo debe describir con detalle los comandos usados para validar. Existe para que el capítulo Validación pueda concentrarse en metodología y resultados, dejando aquí la información operativa. Lo ideal es incluir ejemplos de \texttt{make test-ghc} y \texttt{make test-cabal}.

Para completarla, se debe documentar la estructura de entrada bajo \verb|experiments/| y la estructura de salida bajo \verb|tests/|. También conviene explicar la diferencia entre backend GHC directo y backend Cabal para dependencias externas como \texttt{probability}.
